\documentclass[12pt,a4paper]{article}
\usepackage{fontspec}
\usepackage{polyglossia}
\setmainfont{TeX Gyre Termes}
\newfontfamily\cyrillicfont{TeX Gyre Termes}
\setdefaultlanguage{ukrainian}
\setotherlanguage{english}
\usepackage{amsmath,amssymb,geometry,graphicx}
\geometry{margin=25mm}
\setlength{\parindent}{1.25cm}
\setlength{\parskip}{0.25em}
\begin{document}
\title{Tezy 2024 Оптимізація методу розв'язання СЛАР для полiточкових перетворень}
\maketitle
МЕТОДУ РОЗВ'ЯЗАННЯ СЛАР ДЛЯ ПОЛIТОЧКОВИХ ПЕРЕТВОРЕНЬ Для побудови систем генерації обєктів використовуються політочкові перетворення. Математичний апарат даного методу потребує розвязувати чисельну кількість систем лінійних однорідних рівнянь. Розвязувати системи однорідних лінійних рівнянь можна багатьма способами наприклад методом послідовного виключення (Гаус) та його модифікацією методом Гауса - Жордана. Основна відміність вказаних методів полягає в тому, що метод Гауса приводить систему лінійних однорідних рівнянь до верхньої трикутної матриці (рядкова ступінчаста форма), тоді як метод Гауса-Жордана приводит систему до однорідної (скорочена рядкова ступінчаста форма). Виникає питання, який з цих методів є найбільш ефективни для вирішення задачі методом політочкових перетворень на практиці. Експерементальним способом було проаналізован о, що звичайний метод Гауса на 34\% є ефективнішим за метод Гауса -Жордана. Це обумовлено тим, що метод Гауса -Жордана потребує додаткових операцій нормалізації рядків. Потрібно зазначити, обидва методи можуть бути нестійкими коли працюють з малими числами або мають значні відмінності в порядку між коефіцієнтами. Використання часткового чи повного вибору головного елемента може покращити стійкість системи лінійних однорідних рівнянь для задач, що потребують політочкових перетворень
\par
\end{document}
